\chapter*{Agradecimientos}
\addcontentsline{toc}{chapter}{Agradecimientos}

\ph{En este apartado se incluye una breve nota personal del autor/a en la que se reconoce
el apoyo recibido durante la realización del trabajo. Suele citarse a:}

\begin{itemize}
  \item \ph{Tutor/es académicos: por su orientación, supervisión y disponibilidad.}
  \item \ph{Familiares y/o personas cercanas: por el apoyo personal a lo largo del proceso.}
  \item \ph{Compañeros, grupo de investigación o entidades colaboradoras (si procede).}
\end{itemize}

\ph{Redacta aquí los agradecimientos en primera persona, con un tono cercano pero formal.}

\begin{flushright}
\textit{\ph{Ciudad, mes año}}
\end{flushright}

\vspace{3em}

{\noindent\normalfont\Large\scshape Acknowledgements}\\[-0.4em]
\noindent\rule{\textwidth}{0.4pt}
\addcontentsline{toc}{chapter}{Acknowledgements}

\vspace{1em}

\ph{English version of the acknowledgements (optional, recommended for most universities).
Mirror the same structure as the Spanish version: supervisors, family/close people,
colleagues or collaborating institutions.}

\begin{flushright}
\textit{\ph{City, month year}}
\end{flushright}
